% GENERATED FILE. Edit canonical lessons, then run npm run generate:books.
% canonical-source-hash: fnv1a64:ffc1e06f0a034faf
% canonical-lessons: MW-C286-imtihan, MW-C286-hisab, MW-C286-natijo, MW-C286-darjo, MW-C286-akno

\chapter{Test Day}
\label{ch:mw-test-day}
\hlchaptermodality{286}

\begin{chapteropening}
\textbf{By the end of this chapter:} \emph{I can talk about a test, my class and my result.}

Complete the last lesson of chapter 286: I can talk about a test, my class and my result.
\end{chapteropening}

\section[\texorpdfstring{\mw{इम्तिहान} (imtihān)}{इम्तिहान (imtihān)}]{\mw{इम्तिहान} (imtihān) — an exam}

\label{lesson:MW-C286-imtihan}



\begin{quote}
Before the new one: say the Marwadi for to learn by heart, then the Marwadi for to pronounce.
\end{quote}

\subsection*{You'll want to know: \mw{इम्तिहान}}
\textbf{\mw{इम्तिहान}} — \emph{imtihān} — \textquotedblleft{}an exam\textquotedblright{}.

\begin{quote}
\mw{काल} \mw{इम्तिहान} \mw{है।}
\end{quote}

\begin{quote}
\emph{kāl imtihān hai.}
\end{quote}

\begin{quote}
There's an exam tomorrow.
\end{quote}

A masculine noun.

\begin{cousinweb}
\textbf{\mw{इम्तिहान}} (\emph{imtihān}) is usually traced, through Persian \emph{imtihān}, to Arabic \emph{imtiḥān}, \textquotedblleft{}a test, a trial\textquotedblright{}.
\end{cousinweb}

\begin{grammarlens}[title={What you've built}]
One more word for this chapter.
\end{grammarlens}

\subsection*{Your turn}
\begin{itemize}
\raggedright
  \item \emph{Say it:} \emph{imtihān}
  \item \emph{Say it:} \emph{imtihān}, once more
  \item \emph{Say it:} \emph{kāl imtihān hai.} — There's an exam tomorrow.
  \item \emph{Recall:} say the Marwadi for to learn by heart, then the Marwadi for to pronounce, then say \emph{imtihān} again
\end{itemize}

\begin{tcolorbox}[breakable,colback=teal!4,colframe=teal!35!black,title={Before you move on}]
What does \mw{इम्तिहान} mean? (\textquotedblleft{}An exam\textquotedblright{}.) Say it once more.
\end{tcolorbox}
\section[\texorpdfstring{\mw{हिसाब} (hisāb)}{हिसाब (hisāb)}]{\mw{हिसाब} (hisāb) — sums}

\label{lesson:MW-C286-hisab}



\begin{quote}
Before the new one: say the Marwadi for to pronounce, then the Marwadi for an exam.
\end{quote}

\subsection*{You'll want to know: \mw{हिसाब}}
\textbf{\mw{हिसाब}} — \emph{hisāb} — \textquotedblleft{}sums\textquotedblright{}.

\begin{quote}
\mw{हिसाब} \mw{दोरो} \mw{कोनी।}
\end{quote}

\begin{quote}
\emph{hisāb doro konī.}
\end{quote}

\begin{quote}
Sums aren't hard.
\end{quote}

A masculine noun.

\begin{cousinweb}
\textbf{\mw{हिसाब}} (\emph{hisāb}) is usually traced, through Persian \emph{hisāb}, to Arabic \emph{ḥisāb}, \textquotedblleft{}reckoning\textquotedblright{}: sums at school, and the reckoning of money.
\end{cousinweb}

\begin{grammarlens}[title={What you've built}]
One more word for this chapter.
\end{grammarlens}

\subsection*{Your turn}
\begin{itemize}
\raggedright
  \item \emph{Say it:} \emph{hisāb}
  \item \emph{Say it:} \emph{hisāb}, once more
  \item \emph{Say it:} \emph{hisāb doro konī.} — Sums aren't hard.
  \item \emph{Recall:} say the Marwadi for to pronounce, then the Marwadi for an exam, then say \emph{hisāb} again
\end{itemize}

\begin{tcolorbox}[breakable,colback=teal!4,colframe=teal!35!black,title={Before you move on}]
What does \mw{हिसाब} mean? (\textquotedblleft{}Sums\textquotedblright{}.) Say it once more.
\end{tcolorbox}
\section[\texorpdfstring{\mw{नतीजो} (natījo)}{नतीजो (natījo)}]{\mw{नतीजो} (natījo) — a result}

\label{lesson:MW-C286-natijo}



\begin{quote}
Before the new one: say the Marwadi for an exam, then the Marwadi for sums.
\end{quote}

\subsection*{You'll want to know: \mw{नतीजो}}
\textbf{\mw{नतीजो}} — \emph{natījo} — \textquotedblleft{}a result\textquotedblright{}.

\begin{quote}
\mw{म्हारो} \mw{नतीजो} \mw{चोखो} \mw{है।}
\end{quote}

\begin{quote}
\emph{mhāro natījo chokho hai.}
\end{quote}

\begin{quote}
My result is good.
\end{quote}

A masculine noun.

\begin{cousinweb}
\textbf{\mw{नतीजो}} (\emph{natījo}) is usually traced, through Persian \emph{natīja}, to Arabic \emph{natīja}, \textquotedblleft{}a result, an outcome\textquotedblright{}; Marwadi gives it the -\mw{ो} ending of a masculine noun.
\end{cousinweb}

\begin{grammarlens}[title={What you've built}]
One more word for this chapter.
\end{grammarlens}

\subsection*{Your turn}
\begin{itemize}
\raggedright
  \item \emph{Say it:} \emph{natījo}
  \item \emph{Say it:} \emph{natījo}, once more
  \item \emph{Say it:} \emph{mhāro natījo chokho hai.} — My result is good.
  \item \emph{Recall:} say the Marwadi for an exam, then the Marwadi for sums, then say \emph{natījo} again
\end{itemize}

\begin{tcolorbox}[breakable,colback=teal!4,colframe=teal!35!black,title={Before you move on}]
What does \mw{नतीजो} mean? (\textquotedblleft{}A result\textquotedblright{}.) Say it once more.
\end{tcolorbox}
\section[\texorpdfstring{\mw{दरजो} (darjo)}{दरजो (darjo)}]{\mw{दरजो} (darjo) — a class at school}

\label{lesson:MW-C286-darjo}



\begin{quote}
Before the new one: say the Marwadi for sums, then the Marwadi for a result.
\end{quote}

\subsection*{You'll want to know: \mw{दरजो}}
\textbf{\mw{दरजो}} — \emph{darjo} — \textquotedblleft{}a class at school\textquotedblright{}.

\begin{quote}
\mw{यो} \mw{दरजो} \mw{बड़ो} \mw{है।}
\end{quote}

\begin{quote}
\emph{yo darjo baṛo hai.}
\end{quote}

\begin{quote}
This class is big.
\end{quote}

A masculine noun.

\begin{cousinweb}
\textbf{\mw{दरजो}} (\emph{darjo}) is usually traced, through Persian \emph{darja}, to Arabic \emph{daraja}, \textquotedblleft{}a step, a degree\textquotedblright{}; Marwadi gives it the -\mw{ो} ending of a masculine noun.
\end{cousinweb}

\begin{grammarlens}[title={What you've built}]
One more word for this chapter.
\end{grammarlens}

\subsection*{Your turn}
\begin{itemize}
\raggedright
  \item \emph{Say it:} \emph{darjo}
  \item \emph{Say it:} \emph{darjo}, once more
  \item \emph{Say it:} \emph{yo darjo baṛo hai.} — This class is big.
  \item \emph{Recall:} say the Marwadi for sums, then the Marwadi for a result, then say \emph{darjo} again
\end{itemize}

\begin{tcolorbox}[breakable,colback=teal!4,colframe=teal!35!black,title={Before you move on}]
What does \mw{दरजो} mean? (\textquotedblleft{}A class at school\textquotedblright{}.) Say it once more.
\end{tcolorbox}
\section[\texorpdfstring{\mw{आंकणो} (ā̃kṇo)}{आंकणो (ā̃kṇo)}]{\mw{आंकणो} (ā̃kṇo) — to mark schoolwork; to assess}

\label{lesson:MW-C286-akno}



\begin{quote}
Before the new one: say the Marwadi for a result, then the Marwadi for a class at school.
\end{quote}

\subsection*{You'll want to know: \mw{आंकणो}}
\textbf{\mw{आंकणो}} — \emph{ā̃kṇo} — \textquotedblleft{}to mark schoolwork; to assess\textquotedblright{}.

\begin{quote}
\mw{मास्टर} \mw{हिसाब} \mw{आंकै} \mw{है।}
\end{quote}

\begin{quote}
\emph{māsṭar hisāb ā̃kai hai.}
\end{quote}

\begin{quote}
The teacher marks the sums.
\end{quote}

Past (I …): \mw{म्हैं} \mw{हिसाब} \mw{आंक्यो}, \emph{mhaĩ hisāb ā̃kyo}.

\begin{cousinweb}
\textbf{\mw{आंकणो}} (\emph{ā̃kṇo}) is inherited from Sanskrit \emph{aṅkayati}, \textquotedblleft{}marks, brands\textquotedblright{}; CDIAL gives no Marwari form in this entry, but lists Gujarati \emph{ā̃kvũ} and Hindi \emph{ā̃knā}, both \textquotedblleft{}to mark, appraise\textquotedblright{} (CDIAL 104).
\end{cousinweb}

\begin{grammarlens}[title={What you've built}]
One more word for this chapter.
\end{grammarlens}

\subsection*{Your turn}
\begin{itemize}
\raggedright
  \item \emph{Say it:} \emph{ā̃kṇo}
  \item \emph{Say it:} \emph{ā̃kṇo}, once more
  \item \emph{Say it:} \emph{māsṭar hisāb ā̃kai hai.} — The teacher marks the sums.
  \item \emph{Recall:} say the Marwadi for a result, then the Marwadi for a class at school, then say \emph{ā̃kṇo} again
\end{itemize}

\begin{tcolorbox}[breakable,colback=teal!4,colframe=teal!35!black,title={Before you move on}]
What does \mw{आंकणो} mean? (\textquotedblleft{}To mark schoolwork; to assess\textquotedblright{}.) Say it once more.
\end{tcolorbox}
